\chapter{Melena (Black Stool)}

\section*{Definition}
Melena is black, tarry, often sticky stool caused most often by bleeding in the upper gastrointestinal tract. Blood is digested as it passes through the stomach and small intestine, producing a characteristic dark appearance and strong odor. Iron, bismuth, and some foods can darken stool without bleeding.

\section*{What Happens in Your Body}
Bleeding from the esophagus, stomach, or duodenum is exposed to acid and digestive enzymes. Hemoglobin is altered and mixes with stool, creating black tarry feces. Significant blood loss can reduce circulating volume and oxygen delivery, causing shock, anemia, fainting, or organ injury.

\section*{Causes}
\begin{itemize}
  \item Peptic ulcer disease, gastritis, duodenitis, or bleeding from \textit{H. pylori} infection
  \item Esophageal or gastric varices, especially with advanced liver disease
  \item Mallory-Weiss tear after forceful vomiting or retching
  \item Upper gastrointestinal cancer, vascular lesions, or swallowed blood
  \item Aspirin, NSAIDs, anticoagulants, antiplatelet drugs, or clotting disorders
\end{itemize}

\section*{When to See a Doctor}
Seek emergency care for black tarry stool with vomiting blood or coffee-ground material, fainting, severe weakness, confusion, rapid heartbeat, shortness of breath, pale or clammy skin, or abdominal pain. Any new unexplained melena requires prompt medical assessment, especially when taking blood thinners or NSAIDs.

\section*{Related Symptoms}
\begin{itemize}
  \item Dizziness, fatigue, pallor, weakness, or exercise intolerance from anemia
  \item Upper abdominal pain, nausea, vomiting, or indigestion
  \item Red blood in vomit, low blood pressure, fast pulse, or reduced urine output in major bleeding
\end{itemize}

\section*{Self-Care / Home Management}
Do not take aspirin, ibuprofen, naproxen, or iron supplements for the symptom unless a clinician advises it. Do not stop prescribed anticoagulants without urgent medical guidance. Record stool appearance, medications, and onset, and arrange immediate evaluation rather than waiting for the next bowel movement.

\section*{How Your Doctor Will Evaluate You}
\begin{itemize}
  \item Assessment includes vital signs, circulation, abdominal examination, medication history, and examination for other bleeding.
  \item Blood count, urea and electrolytes, coagulation studies, liver tests, blood typing, and crossmatching may be required.
  \item Upper endoscopy usually identifies and may treat the bleeding source; CT angiography or angiography is used in selected cases.
\end{itemize}

\section*{Treatment Options}
Treatment may include intravenous fluids, blood products when indicated, acid suppression, reversal of anticoagulation when appropriate, endoscopic hemostasis, variceal therapy, interventional radiology, or surgery. Treatment also addresses \textit{H. pylori}, medication injury, liver disease, or malignancy.

\section*{References}
\begin{thebibliography}{99}
\bibitem{melena:ref1} Laine L, Barkun AN, Saltzman JR, et al. ACG Clinical Guideline: Upper Gastrointestinal and Ulcer Bleeding. \textit{Am J Gastroenterol}. 2021;116:899--917.
\bibitem{melena:ref2} Gralnek IM, et al. Endoscopic diagnosis and management of nonvariceal upper gastrointestinal hemorrhage. \textit{Endoscopy}. 2021;53:300--332.
\bibitem{melena:ref3} Garcia-Tsao G, et al. Portal hypertensive bleeding in cirrhosis. \textit{Hepatology}. 2016;63:472--492.
\end{thebibliography}
